%%%%%%%%%%%%%%%%%%%%%%%%%%%%%%%%%%%%%%%%%
% eBook 
% LaTeX Template
% Version 1.0 (29/12/14)
%
% This template has been downloaded from:
% http://www.LaTeXTemplates.com
%
% Original author:
% Luis Cobo (luiscobogutierrez@gmail.com) with extensive modifications by:
% Vel (vel@latextemplates.com)
%
% License:
% CC BY-NC-SA 3.0 (http://creativecommons.org/licenses/by-nc-sa/3.0/)
%
%%%%%%%%%%%%%%%%%%%%%%%%%%%%%%%%%%%%%%%%%

%----------------------------------------------------------------------------------------
%	DOCUMENT CONFIGURATIONS AND INFORMATION
%----------------------------------------------------------------------------------------

\documentclass[oneside,11pt]{memoir} % Font size

\input{structure.tex} % Include the file that specifies the document structure and layout

\title{ॐ | वैदिक गणित | विचार संग्रह } % Book title
\author{Amit Kumar Singh} % Author
\newcommand{\edition}{Apr-2020} % Book edition

%----------------------------------------------------------------------------------------

\begin{document}

%----------------------------------------------------------------------------------------
%	TITLE PAGE
%----------------------------------------------------------------------------------------

\thispagestyle{empty} % Suppress page numbering
\ThisCenterWallPaper{1}{_images/FATHER-OF-VEDIC-MATHS.jpg} % Add the background image, the first argument is the scaling - adjust this as necessary so the image fits the entire page

ॐ | वैदिक गणित | विचार संग्रह 
% \begin{tikzpicture}[remember picture,overlay]
% \node [rectangle, rounded corners, fill=white, opacity=0.75, anchor=south west, minimum width=3cm, minimum height=6cm] (box) at (-0.5,-10) (box){}; % White rectangle - "minimum width/height" adjust the width and height of the box; "(-0.5,-10)" adjusts the position on the page
% \node[anchor=west, color01, xshift=-1.5cm, yshift=-0.4cm, text width=2.9cm, font=\sffamily\scriptsize] at (box.north){\edition}; % "Text width" adjusts the wrapping width, "xshift/yshift" adjust the position relative to the white rectangle
% \node[anchor=west, color01, xshift=-1.5cm, yshift=-2cm, text width=2.9cm, font=\sffamily\bfseries\scshape\Large] at (box.north){\thetitle}; % "Text width" adjusts the wrapping width, "xshift/yshift" adjust the position relative to the white rectangle
% \node[anchor=west, color01, xshift=-1.5cm, yshift=-5cm, text width=2.9cm, font=\sffamily\bfseries] at (box.north){\theauthor}; % "Text width" adjusts the wrapping width, "xshift/yshift" adjust the position relative to the white rectangle
% \end{tikzpicture}

\newpage % Make sure the following content is on a new page

%----------------------------------------------------------------------------------------
%	TABLE OF CONTENTS
%----------------------------------------------------------------------------------------

\tableofcontents % Prints the table of contents

%----------------------------------------------------------------------------------------
%	INTRODUCTION SECTION
%----------------------------------------------------------------------------------------

\chapter*{परिचय} % Introduction chapter suppressed from the table of contents

\begin{quote}

\noindent\rule{7.5cm}{0.4pt}

ॐ गं गणपतये नमः 

\noindent\rule{7.5cm}{0.4pt}

'शारदा शारदाभौम्वदना। वदनाम्बुजे।

सर्वदा सर्वदास्माकमं सन्निधिमं सन्निधिमं क्रिया तू।'

श्रीं ह्रीं सरस्वत्यै स्वाहा।

ॐ ह्रीं ऐं ह्रीं सरस्वत्यै नमः।

\noindent\rule{7.5cm}{0.4pt}


\end{quote}

\section{यह पुस्तक क्यों?}

 
\footnote{}.

%----------------------------------------------------------------------------------------
%	BOOK PART
%----------------------------------------------------------------------------------------

\part{वैदिक गणित का सार अपने शब्दों में। }

%----------------------------------------------------------------------------------------
%	CHAPTER ZERO 
%----------------------------------------------------------------------------------------


\chapter{रेखांक?}

\section{संख्याएँ (Numbers)}
शून्य (०), एक से नौ (१-९), रेखांक एक से रेखांक नौ ($\bar{१}$-$\bar{९}$)
\section{आधार (base)}
एक संख्या का आधार बताता है की संख्या में कितने अलग चिन्ह होंगे। द्वी अंकीय संख्या का आधार २ होता है। द्वीअंकीय संख्या के दो चिन्ह ० और १ होता है. दशमलव संख्या में दस चिन्ह होते हैं। दसमलव संख्या के चिन्ह ०, १,  २, ३, ४,५,  ६, ७, ८,९ मान्य है। आधार को ठीक से समझने के लिए हम यह जान लेते हैं की आधार का उपयोग क्या है। यदि हम दस की संख्या दर्शाना चाहते हैं तो इस तरह दर्शाएंगे। 

\begin{enumerate}
\item{द्विअंकीय प्रणाली:}
$१ * (२)^{३} + ० * (२)^{२} + १ * (२)^{१} + १ * (२)^{०}$ = $१*८+०+२+०=१०$ = $८+२=१०$
\item{दशमलव प्रणाली:}
$१ * (१०)^{१} + ० * (१०)^{०}$ = $१०+०=१०$
\end{enumerate}

\section{आधार के पास गुणा}

\begin{tikzpicture}

\draw[step=1cm,gray,very thin] (0,0) grid (5,5);

\node (103) at (1,3) [number] {103};


\node (89) [number, below of=103] {89};





\end{tikzpicture}

\begin{tikzpicture}

\draw[step=1cm,gray,very thin] (0,0) grid (5,5);

\node (103) at (1,3) [number] {103};

\node (100) at (2,4) [number] {\underline{100}};
\node (aadhar) at (4,4) [number] {आधार};
\node (89) [number, below of=103] {89};





\end{tikzpicture}

\begin{tikzpicture}

\draw[step=1cm,gray,very thin] (0,0) grid (5,5);

\node (103) at (1,3) [number] {103};
\node (3) [number, right of=103] {3};
\node (poorak) [number, right of=3] {पूरक};
\node (100) [number, above of=3] {\underline{100}};
\node (89) [number, below of=103] {89};
\node (-11) [-number, right of=89] {$\bar{11}$};
\node (poorak2) [-number, right of=-11] {पूरक};



\end{tikzpicture}

\begin{tikzpicture}

\draw[step=1cm,gray,very thin] (0,0) grid (5,5);

\node (103) at (1,3) [number] {103};
\node (3) [number, right of=103] {3};
\node (100) [number, above of=3] {\underline{100}};
\node (89) [number, below of=103] {89};
\node (-11) [-number, right of=89] {$\bar{11}$};
\node (-33) [-number, below of=-11] {$\bar{33}$};
\node (guna) [-number, right of=-33, xshift=1cm] {गुणा पूरक };



\end{tikzpicture}


\begin{tikzpicture}

\draw[step=1cm,gray,very thin] (0,0) grid (5,5);

\node (103) at (1,3) [number] {103};
\node (3) [number, right of=103] {3};
\node (100) [number, above of=3] {\underline{100}};
\node (89) [number, below of=103] {89};
\node (-11) [-number, right of=89] {$\bar{11}$};
\node (-33) [-number, below of=-11] {$\bar{33}$};
\node (92) [number, below of=89] {92};
\node (tirabhyam) [number, right of=-33, xshift=1cm] {तिरभयाम};


\end{tikzpicture}


\begin{tikzpicture}

\draw[step=1cm,gray,very thin] (0,0) grid (5,5);

\node (103) at (1,3) [number] {103};
\node (3) [number, right of=103] {3};
\node (100) [number, above of=3] {\underline{100}};
\node (89) [number, below of=103] {89};
\node (-11) [-number, right of=89] {$\bar{11}$};
\node (-33) [-number, below of=-11] {$\bar{33}$};
\node (92) [number, below of=89] {92};
\node (67) at (4,0) [number] {67};
\node (nikhilam) at (4,2) [number] {निखिलं};


\end{tikzpicture}


\begin{tikzpicture}

\draw[step=1cm,gray,very thin] (0,0) grid (5,5);

\node (103) at (1,3) [number] {103};
\node (3) [number, right of=103] {3};
\node (100) [number, above of=3] {\underline{100}};
\node (89) [number, below of=103] {89};
\node (-11) [-number, right of=89] {$\bar{11}$};
\node (-33) [-number, below of=-11] {$\bar{33}$};
\node (92) [number, below of=89] {92};
\node (67) at (4,0) [number] {67};
\node (91) at (3,0) [number] {91};
\node (eknupu) at (4,2) [number] {एकन्यूनेन};


\end{tikzpicture}

%----------------------------------------------------------------------------------------
%	CHAPTER ONE
%----------------------------------------------------------------------------------------


\chapter{Sootra is seed | sow selectively}

My serious interest in thoughts of my ancestors began with संस्कृत. I was a teenager when my grandfather left his body. I remember the rituals performed by brahmins of my village beside one of its lakes. The sounds eminating from the chants performed by the brahmins was mesmerizing and it stayed in me like a seed waiting for right season. ``Where do I come from?'' is the question I began asking there on.

I am a Rajpoot. That segment among Rajpoots class of society that rule over 36 other classes.

%----------------------------------------------------------------------------------------
%	CHAPTER TWO
%----------------------------------------------------------------------------------------

\chapter{Rituals to sustain | take responsibilities}


%----------------------------------------------------------------------------------------
%	CHAPTER THREE
%----------------------------------------------------------------------------------------

\chapter{Give away fruits | to who deserves}



%----------------------------------------------------------------------------------------

\end{document}